\begin{statementbox}
\textbf{条件。} $a,b>0$，直线 $l:Ax+By+C=0$，$A^2+B^2>0$。
交点用 $P(x_1,y_1),Q(x_2,y_2)$ 表示，避免与系数 $A,B$ 混用。
\textbf{消去 $y$ 必须有 $B\ne0$；消去 $x$ 必须有 $A\ne0$。}

\textbf{一、横开口：$\dfrac{x^2}{a^2}-\dfrac{y^2}{b^2}=1$。}
记 $M=B^2b^2-A^2a^2$，$H=C^2+M$。
在各自消元条件满足时，得到
\[
 \begin{aligned}
 Mx^2-2ACa^2x-a^2(C^2+B^2b^2)&=0,\\
 My^2+2BCb^2y+b^2(C^2-A^2a^2)&=0.
 \end{aligned}
\]
当 $M\ne0$，韦达关系为
\[
 \boxed{x_1+x_2=\frac{2ACa^2}{M}},\qquad
 \boxed{x_1x_2=-\frac{a^2(C^2+B^2b^2)}{M}},
\]
\[
 \boxed{y_1+y_2=-\frac{2BCb^2}{M}},\qquad
 \boxed{y_1y_2=\frac{b^2(C^2-A^2a^2)}{M}}.
\]
\textbf{交点分类：} $M\ne0$ 时，$H>0$ 为两个不同实交点，$H=0$ 为相切，$H<0$ 为无实交点。此时
\[
 \Delta_x=4a^2b^2B^2H\ (B\ne0),\qquad
 \Delta_y=4a^2b^2A^2H\ (A\ne0).
\]
实交点的和、积要求 $H\ge0$；相切时重计切点。

\textbf{退化必须另算。} $M=0$ 时 $A,B$ 均非零，直线平行于一条渐近线：
\[
 \begin{cases}
 C\ne0: &\text{恰有一个有限交点，且不是切线；}\\
 C=0: &\text{直线就是渐近线，没有有限交点。}
 \end{cases}
\]
不能把 $M=0$ 代入韦达分式，也不能用二次方程的判别式分类。
\end{statementbox}
\clearpage
\begin{statementbox}
\textbf{二、纵开口：$\dfrac{y^2}{a^2}-\dfrac{x^2}{b^2}=1$。}
记 $N=A^2b^2-B^2a^2$，$K=C^2+N$。
消去 $y$ 仍需 $B\ne0$，消去 $x$ 仍需 $A\ne0$；相应方程为
\[
 \begin{aligned}
 Nx^2+2ACb^2x+b^2(C^2-B^2a^2)&=0,\\
 Ny^2-2BCa^2y-a^2(C^2+A^2b^2)&=0.
 \end{aligned}
\]
当 $N\ne0$，得到
\[
 \boxed{x_1+x_2=-\frac{2ACb^2}{N}},\qquad
 \boxed{x_1x_2=\frac{b^2(C^2-B^2a^2)}{N}},
\]
\[
 \boxed{y_1+y_2=\frac{2BCa^2}{N}},\qquad
 \boxed{y_1y_2=-\frac{a^2(C^2+A^2b^2)}{N}}.
\]
位置关系将前页的 $(M,H)$ 换为 $(N,K)$ 即可：$N\ne0$ 时按 $K$ 的正、零、负分别得到两交点、相切、无交点。
$N=0$ 时，$C\ne0$ 恰有一个有限交点，$C=0$ 则是渐近线、无有限交点。

\textbf{三、中点与弦长（适用于水平、竖直及斜弦）。}
当存在两个不同实交点时，弦中点 $S$ 与弦长为
\[
 \begin{array}{c|c|c}
 & S & PQ\\[3pt]\hline
 \text{横开口} & \left(\dfrac{ACa^2}{M},-\dfrac{BCb^2}{M}\right)
 & \dfrac{2ab\sqrt{(A^2+B^2)H}}{|M|}\\[10pt]
 \text{纵开口} & \left(-\dfrac{ACb^2}{N},\dfrac{BCa^2}{N}\right)
 & \dfrac{2ab\sqrt{(A^2+B^2)K}}{|N|}
 \end{array}
\]
横开口须 $M\ne0,H>0$；纵开口须 $N\ne0,K>0$。
分母要取\textbf{绝对值}，因为二次项系数可以为负。
相切时中点公式给出切点，长度公式给出 $0$，不再是非零弦。
\end{statementbox}
\clearpage
